% Part: lambda-calculus
% Chapter: introduction
% Section: church-rosser

\documentclass[../../../include/open-logic-section]{subfiles}

\begin{document}

\olfileid[tr]{lam}{int}{cr}
\olsection{Church--Rosser Özelliği}

\begin{thm}
\ollabel{thm:church-rosser}
$M$, $N_1$ ve $N_2$, $M\red N_1$ ve $M\red N_2$ olacak biçimde terimler
olsun. O zaman $N_1\red P$ ve $N_2\red P$ olacak biçimde bir $P$ terimi
vardır.
\end{thm}

\begin{cor}
$M$'nin normal biçime indirgenebildiğini varsayın. O zaman bu normal biçim
tektir.
\end{cor}

\begin{proof}
$M\red N_1$ ve $M\red N_2$ ise önceki teorem gereği $N_1$ ile $N_2$'nin
ikisi de bir $P$ terimine indirgenir. $N_1$ ile $N_2$ normal biçimdeyse bu,
ancak $N_1\ident P\ident N_2$ olduğunda mümkündür.
\end{proof}

Son olarak iki $M$ ve $N$ terimi ortak bir terime indirgeniyorsa
\emph{$\beta$-eşdeğer}, kısaca \emph{eşdeğer} diyeceğiz; başka bir deyişle
$M\red P$ ve $N\red P$ olacak biçimde bir $P$ varsa. Bu $M\equal[\beta]N$
ile yazılır. \olref{thm:church-rosser} kullanılarak $\equal[\beta]$'nın bir
denklik bağıntısı olduğu ve ayrıca her $M$ ve $N$ için $M\red N$ ya da
$N\red M$ ise $M\equal[\beta]N$ olduğu denetlenebilir. (Hatta
$\equal[\beta]$'nın bu özelliği taşıyan \emph{en küçük} denklik bağıntısı
olduğu gösterilebilir.)

\end{document}

